\documentclass[11pt]{article}
\usepackage[margin=1in]{geometry}
\usepackage{amsmath,amssymb,amsthm,mathtools}
\usepackage[colorlinks=true,linkcolor=blue,citecolor=blue,urlcolor=blue]{hyperref}
\hypersetup{pdftitle={Horn kernels and normalization of diagonal simplicial tensor modules},pdfauthor={Florian Lengyel}}
\newtheorem{theorem}{Theorem}[section]
\newtheorem{proposition}[theorem]{Proposition}
\newtheorem{corollary}[theorem]{Corollary}
\newtheorem{lemma}[theorem]{Lemma}
\theoremstyle{definition}
\newtheorem{definition}[theorem]{Definition}
\theoremstyle{remark}
\newtheorem{remark}[theorem]{Remark}
\newtheorem{example}[theorem]{Example}
\DeclareMathOperator{\im}{im}
\DeclareMathOperator{\rank}{rank}
\DeclareMathOperator{\Stab}{Stab}
\newcommand{\SII}[2]{\genfrac{\{}{\}}{0pt}{}{#1}{#2}}
\title{Horn kernels and normalization of diagonal simplicial tensor modules}
\author{Florian Lengyel\\\small The City University of New York}
\date{}
\begin{document}
\maketitle
\begin{abstract}
Deleting or repeating the same index in every axis of a tensor gives
concrete examples of simplicial modules over a nonzero commutative ring.
We use these examples to study horn filling, normalization, chain
contractions, and homology through coordinate calculations, supported by
a Python implementation. Horn kernels have explicit bases, and fillers
are unique precisely in dimensions above the number of tensor axes.
The normalized complex splits into two-term complexes with identity
differential, giving explicit equivariant contractions. Two shapes give
isomorphic simplicial modules exactly when they differ by a permutation
and addition of the same integer to every entry. Although the full module
is contractible, a submodule generated by a single matrix can have
nonzero first homology. We describe generated submodules as marked quotients. An appendix
studies their kernel families over an infinite field, including an
explicit parameterization for $3\times3$ matrices.
\end{abstract}

\section{Introduction}
Matrices and higher-order arrays provide a computational way to study
simplicial modules. Deleting or repeating the same index in every axis
defines face and degeneracy operators; allowing the dimensions to vary
together gives a simplicial module. The \emph{SimplicialTensors}
repository \cite{SimplicialTensors} implements these operations, the
alternating boundary, and horn filling in Python. This paper develops
the tensor example as a concrete route through the corresponding
constructions of simplicial algebra.

For instance, horn equations ask whether a compatible collection of
diagonal faces determines a tensor and, when it does not, how its
fillers differ. Normalization removes the chains obtained by degeneracy
operators. A chain contraction then supplies an explicit filling of
each cycle in the full tensor module. These constructions can be
examined both on coordinate tensors and in array computations.

These questions admit coordinate answers. A face kills a coordinate
tensor precisely when its index tuple contains the deleted index.
Consequently a horn kernel has a basis indexed by tuples covering the
prescribed face indices. Its rank follows by inclusion--exclusion, and
horn fillers are unique exactly above the number of tensor axes
(Theorem~\ref{thm:hornbasis}). The same basis splits the normalized
complex into disks (Proposition~\ref{prop:disks}), yielding explicit
contractions. Two shapes give isomorphic simplicial modules exactly when
they differ by a permutation and addition of the same integer to every entry.

The ambient simplicial module is contractible, but a submodule generated
by one tensor need not be. A $2\times2$ matrix gives an elementary example
with one-dimensional first homology (Example~\ref{ex:matrix-cycle}).
We describe a generated submodule as a marked
quotient of the free simplicial module on a simplex. Normalization makes
its kernel finite-dimensional data (Section~\ref{sec:generated}).
Appendix~\ref{app:kernel-families} proves that the closure of the generic
kernel image is rational over an infinite field and computes the image
and quotient homology for shape $(3,3)$. This appendix additionally uses
Grassmannians, incidence schemes, and function fields.

We use the normalization and Dold--Kan theorems
\cite{Dold1958,Kan1958,Weibel}. The horn-space isomorphism of
\cite[Lemma 2.10]{RonchiZhu}, stated there for simplicial objects in
finite-dimensional real vector spaces, follows here from a
horn-complement argument over $A$. All coordinate calculations are valid over
a nonzero commutative ring; Section~\ref{sec:generated} works over a field.

\section{The construction}
Throughout, $A$ is a nonzero commutative ring with identity, $k\geq1$, and
$[p]=\{0,\ldots,p\}$. Let $\Delta$ denote the category of finite nonempty
ordinals and nondecreasing maps. Fix a shape
$\vec s=(n_1,\ldots,n_k)\in\mathbb Z_{>0}^k$ and put
\[
 n=\min_a n_a-1,\qquad c_a=n_a-1-n,\qquad
 I_p=\prod_{a=1}^k[p+c_a].
\]
The integers $c_a=n_a-\min_b n_b$ are the \emph{shape offsets}.
Thus $c_a\geq0$ and at least one $c_a$ is zero. The integer $n$ records
the degree in which tensors of the original shape occur; the offsets
$(c_a)$ determine the entire simplicial module.

For $c\geq0$, put $J_c([p])=[p+c]$. For $\alpha:[p]\to[q]$, define
\[
 J_c(\alpha)(r)=
 \begin{cases}
 \alpha(r),&0\leq r\leq p,\\
 q+t,&r=p+t,\quad1\leq t\leq c.
 \end{cases}
\]
This defines a covariant functor $J_c:\Delta\to\mathbf{Set}$, since
composition acts separately on the initial ordinal and the final $c$
elements. Consequently $I=\prod_a J_{c_a}$ is cosimplicial and
\[
 X_p=X_p(\vec s;A)=A^{I_p},\qquad
 X(\alpha)(T)=T\circ I(\alpha)
\]
defines a functor $X:\Delta^{\mathrm{op}}\to\mathbf{Mod}_A$.
We call $X$ the \emph{diagonal simplicial tensor module} of shape $\vec s$.
Write $E_m$ for the indicator function of $m\in I_p$.

For $p\geq1$ and $0\leq i\leq p$, the index coface
$\Delta_i:I_{p-1}\to I_p$ inserts a gap at $i$ in
every coordinate; its image consists of the tuples avoiding $i$.
For $p\geq0$ and $0\leq i\leq p$, the codegeneracy
$\Sigma_i:I_{p+1}\to I_p$ identifies $i$ and $i+1$ in every coordinate. Thus $d_iT=T\circ\Delta_i$ and
$s_iT=T\circ\Sigma_i$. In particular,
\begin{equation}\label{eq:face}
 d_iE_m=
 \begin{cases}
 0,&i\in\im(m),\\
 E_{\Delta_i^{-1}m},&i\notin\im(m),
 \end{cases}
 \qquad \im(m)=\{m_1,\ldots,m_k\}.
\end{equation}
The nonzero images in this formula are distinct basis elements.
There are no outgoing face maps in degree zero. The associated chain
complex has $X_{-1}=0$ and differential
\[
 \partial_p=\sum_{i=0}^p(-1)^id_i\quad(p\geq1),\qquad \partial_0=0.
\]
For any simplicial module $Y$, we write
$H_p(Y)=\ker\partial_p/\im\partial_{p+1}$.
When $k=1$, $X_p=A^{p+1}$; a face of a two-entry vector is a
one-entry vector.

In the repository's \path{simplicial_tensors.tensor_ops} module,
\texttt{face(T,i)} and \texttt{degen(T,i)} compute these face and
degeneracy maps, and \texttt{bdry(T)} computes the alternating face
boundary in positive degrees. In degree zero, the empty array returned
by \texttt{bdry(T)} represents the zero boundary.
Install with \texttt{pip install -e .}, then run
\texttt{python} \mbox{\texttt{experiments/dstm\_examples.py}} to reproduce
Examples~\ref{ex:matrix-cycle} and~\ref{ex:matrix-horn}. The script uses
integer arrays and exact ranks over $\mathbb Q$: it reports first
homology dimension $1$ for the generated matrix example and $2$ free
parameters for the horn. These calculations illustrate the examples;
they do not implement a general homology algorithm.

The construction is the diagonal of the $k$-fold simplicial module
$A^{\prod_a[p_a+c_a]}$; compare the diagonal construction in
\cite{Quillen1966}. Equivalently,
$X\cong\bigotimes_{a=1}^k A^{J_{c_a}}$, with degreewise tensor product.
All structure maps are defined over $\mathbb Z$, and coordinate bases give
\begin{equation}\label{eq:basechange}
 X(\vec s;A)\otimes_A B\cong X(\vec s;B)
\end{equation}
for every homomorphism $A\to B$.

For $T\in X_n$, its generated simplicial submodule has degree-$p$ part
\[
 \langle T\rangle_p=
 \operatorname{span}_A\{X(\alpha)(T):\alpha:[p]\to[n]\}.
\]
This is closed under faces and degeneracies by functoriality.

\begin{example}\label{ex:matrix-cycle}
Over a field $K$, take shape $(2,2)$, write $E_{ij}=E_{(i,j)}$ and let
\[
 T=E_{01}=\begin{pmatrix}0&1\\0&0\end{pmatrix},
 \qquad Y=\langle T\rangle\subseteq X.
\]
Both faces vanish: $d_0T=d_1T=0$. Consequently $Y_0=0$,
$Y_1=KT$, and
\[
 Y_2=\operatorname{span}_K\{s_0T,s_1T\},\qquad
 s_0T=E_{02}+E_{12},\quad s_1T=E_{01}+E_{02}.
\]
The simplicial identities give
\[
 \partial_2s_0T=T-T+s_0d_1T=0,\qquad
 \partial_2s_1T=s_0d_0T-T+T=0.
\]
Thus $T$ is a nonzero cycle that is not a boundary in $Y$, and
$H_1(Y)\cong K$. Every simplex of $Y$ above degree $1$ is a sum of
degeneracies.

In the ambient complex, however,
\[
 \partial_2E_{12}=T.
\]
This primitive does not belong to $Y_2$. Indeed, expressing $E_{12}$
as $a(E_{02}+E_{12})+b(E_{01}+E_{02})$ would force $b=0$, $a+b=0$,
and $a=1$, a contradiction. Thus a cycle in a simplicial submodule
can bound in the ambient module without bounding in the submodule.
\end{example}

\section{Horn kernels and normalization}
For any simplicial $A$-module $Y$, put
\[
 N_p(Y)=\bigcap_{i=1}^p\ker d_i,\qquad
 D_p(Y)=\sum_{i=0}^{p-1}\im s_i.
\]
Here $N_0(Y)=Y_0$, $D_0(Y)=0$, and the normalized differential is
$d_0$ in positive degrees and zero in degree zero; this is the convention
in \cite[Exercise 8.3.4]{Weibel}. The normalization theorem gives
\mbox{$Y_p=N_p(Y)\oplus D_p(Y)$} \cite[Corollary 1.12]{Dold1958}.
The Dold--Kan theorem gives an equivalence with nonnegative chain
complexes \cite[Theorem 8.4.1]{Weibel}.

For $p\geq1$ and $0\leq j\leq p$, let
\[
 R_{p,j}(Y)=\bigcap_{i\ne j}\ker d_i.
\]
The module $\Lambda_{p,j}(Y)$ of compatible horns consists of tuples
$(y_i)_{i\ne j}$ with $y_i\in Y_{p-1}$ satisfying $d_i y_\ell=d_{\ell-1}y_i$ for
$i<\ell$, $i,\ell\ne j$; for $p=1$ there are no compatibility
conditions. The horn map is
$\Phi_j:Y_p\to\Lambda_{p,j}(Y)$, $y\mapsto(d_i y)_{i\ne j}$.
Every simplicial module is Kan, so $\Phi_j$ is surjective
\cite[Lemma 8.2.8]{Weibel}. Its fibers are torsors for $R_{p,j}(Y)$.

\begin{lemma}\label{lem:hornsplit}
For every simplicial $A$-module $Y$, $p\geq1$ and $0\leq j\leq p$,
\[
 Y_p=R_{p,j}(Y)\oplus D_p(Y).
\]
The restriction $\Phi_j|_{D_p(Y)}$ is an isomorphism. Thus every horn
has a unique filler in the submodule $D_p(Y)$.
\end{lemma}
\begin{proof}
For $0\leq j<p$, the simplicial identities give inverse maps
\[
 F_j=1-s_jd_j:R_{p,j}(Y)\longrightarrow R_{p,j+1}(Y),\qquad
 G_j=1-s_jd_{j+1}:R_{p,j+1}(Y)\longrightarrow R_{p,j}(Y).
\]
Indeed, the required faces vanish by $d_i d_j=d_{j-1}d_i$ for
$i<j$ and $d_js_j=d_{j+1}s_j=1$. For $x\in R_{p,j}$,
$d_{j+1}F_jx=-d_jx$, whence $G_jF_jx=x$; the other composition
is identical with $j$ and $j+1$ interchanged.
Both maps induce the identity modulo $D_p(Y)$. Since
$R_{p,0}(Y)=N_p(Y)$ maps isomorphically onto $Y_p/D_p(Y)$,
the same holds for every $R_{p,j}(Y)$. This proves the direct sum.
The final assertion follows from surjectivity of $\Phi_j$ and
$\ker\Phi_j=R_{p,j}(Y)$.
\end{proof}

For simplicial objects in finite-dimensional real vector spaces the horn-space
isomorphism is also proved in \cite[Lemma 2.10]{RonchiZhu};
Lemma~\ref{lem:hornsplit} gives an argument over arbitrary $A$.

\begin{theorem}\label{thm:hornbasis}
For the diagonal simplicial tensor module $X$, let
$F_j=[p]\setminus\{j\}$. Then
\begin{equation}\label{eq:hornbasis}
 R_{p,j}(X)=\bigoplus_{\substack{m\in I_p\\F_j\subseteq\im(m)}} AE_m,
 \qquad
 \rank R_{p,j}(X)=
 \sum_{t=0}^p(-1)^t\binom pt\prod_{a=1}^k(p+c_a+1-t).
\end{equation}
In particular, $R_{p,j}(X)\ne0$ if and only if $p\leq k$.
Thus horn fillers are nonunique for $1\leq p\leq k$ and unique
for $p>k$, independently of $j$ and of the offsets.
\end{theorem}
\begin{proof}
By \eqref{eq:face}, $d_iT=0$ exactly when $T$ is supported on
tuples containing $i$. Intersecting these support conditions gives
\eqref{eq:hornbasis}. Inclusion--exclusion counts tuples containing
all $p$ elements of $F_j$: forbidding a specified $t$ of these
elements leaves $p+c_a+1-t$ choices in coordinate $a$.
A tuple of length $k$ cannot contain $p>k$ distinct values. If
$p\leq k$, assign the $p$ required values to distinct coordinates
and fill the remaining coordinates arbitrarily in $[p]$.
This gives a basis element, nonzero because $A\ne0$.
\end{proof}


\begin{example}\label{ex:matrix-horn}
For shape $(3,3)$, consider the $2$-horn with missing face $0$ and
prescribed faces
\[
 d_1T=\begin{pmatrix}1&3\\7&9\end{pmatrix},\qquad
 d_2T=\begin{pmatrix}1&2\\4&5\end{pmatrix}.
\]
The compatibility condition $d_1d_2T=d_1d_1T$ holds: both sides
are the one-entry matrix $(1)$. Every filler is of the form
\[
 T=\begin{pmatrix}1&2&3\\4&5&u\\7&v&9\end{pmatrix},
 \qquad u,v\in A.
\]
Thus the difference of two fillers lies in
$R_{2,0}(X)=AE_{(1,2)}\oplus AE_{(2,1)}$, whose basis consists
of the two index pairs containing both $1$ and $2$.
\end{example}

An algebraic $r$-hypergroupoid is a simplicial module whose horns
have unique fillers in every dimension greater than $r$
\cite{Duskin1979,Glenn1982}. Theorem~\ref{thm:hornbasis} says that
$X$ is such an object exactly when $r\geq k$; its least possible
value of $r$ is $k$. In particular, for the original tensor degree
$n=\min_a n_a-1$, the $n$-hypergroupoid condition is $k\leq n$,
with nonunique $n$-horn fillers when $k=n\geq1$.

\section{The normalized complex}
Set $Z_p=\ker(d_0:N_p(X)\to N_{p-1}(X))$ for $p\geq1$ and
$Z_0=X_0$. Since some $c_a=0$, every index in $I_0$ contains zero.
Therefore, in all degrees $p\geq0$,
\begin{equation}\label{eq:cycles}
 Z_p=\bigoplus_{\substack{m\in I_p\\{}[p]\subseteq\im(m)}} AE_m.
\end{equation}

\begin{proposition}\label{prop:disks}
For $p\geq1$ and $0\leq j\leq p$, define
$\iota_j:Z_{p-1}\to R_{p,j}(X)$ by
$\iota_j(E_m)=E_{\Delta_jm}$. Then
\[
 R_{p,j}(X)=Z_p\oplus\iota_jZ_{p-1},\qquad d_j\iota_j=1.
\]
Consequently, writing $D^r(M)$ for the chain complex $M\xrightarrow{1}M$
in degrees $r,r-1$, there is a coordinate-defined isomorphism
\begin{equation}\label{eq:disks}
 N_\bullet(X)\cong\bigoplus_{q=0}^{k-1}D^{q+1}(Z_q).
\end{equation}
In particular $N_\bullet(X)$ is split exact.
\end{proposition}
\begin{proof}
Partition the basis in \eqref{eq:hornbasis} according as the tuple
contains $j$ or avoids $j$. The first part is precisely the basis
of $Z_p$. Deleting the gap at $j$ identifies the second part
bijectively with the basis of $Z_{p-1}$, and \eqref{eq:face}
describes $d_j$ on both parts. Taking $j=0$ pairs each basis
element of $Z_q$ with its coordinatewise shift by $1$ in degree
$q+1$. There are no other basis elements of $N_\bullet(X)$.
Finally $Z_q=0$ for $q\geq k$, because $k$ coordinates cannot
contain $q+1$ distinct values.
\end{proof}

Explicitly, the contraction $h_p:N_p(X)\to N_{p+1}(X)$ is
\[
 h_p(E_m)=
 \begin{cases}
 E_{m+\mathbf1},&0\in\im(m),\\
 0,&0\notin\im(m).
 \end{cases}
\]
Thus $d_0h+hd_0=1$, with the differential interpreted as zero
in degree zero. The bases in \eqref{eq:hornbasis} and
\eqref{eq:cycles} also show that horn kernels, normalized chains
and normalized cycles commute with arbitrary base change.

\begin{corollary}\label{cor:ranks}
Put $b_p=\rank Z_p$ and $b_{-1}=0$. Then
\[
 b_p=\sum_{t=0}^{p+1}(-1)^t\binom{p+1}{t}
          \prod_{a=1}^k(p+c_a+1-t),\qquad
 \rank N_p(X)=b_p+b_{p-1}.
\]
For constant shapes, $c_a=0$ for every $a$, and
\[
 b_p=(p+1)!\SII{k}{p+1},\qquad
 \rank N_p(X)=p!\SII{k}{p}+(p+1)!\SII{k}{p+1},
\]
where $\SII{k}{r}$ is a Stirling number of the second kind.
\end{corollary}
\begin{proof}
Apply inclusion--exclusion to \eqref{eq:cycles} and then use
Proposition~\ref{prop:disks}. In the constant-shape case the
indices covering $[p]$ are exactly the surjections from a $k$-element
set onto a $(p+1)$-element set, counted by
$(p+1)!\SII{k}{p+1}$ \cite[\S1.4]{StanleyEC1}.
\end{proof}

\begin{example}
For $k=2$ and any constant shape $(N,N)$, one has
$X_p=A^{(p+1)\times(p+1)}$ and
\[
 N_\bullet(X)\cong D^1(A)\oplus D^2(A^2).
\]
The normalized ranks are $1,3,2$ in degrees $0,1,2$ and zero
thereafter. In degree $2$, the two basis elements of $R_{2,0}$
are $E_{(1,2)}$ and $E_{(2,1)}$. They map under $d_0$ to
$E_{(0,1)}$ and $E_{(1,0)}$, respectively.
\end{example}

\begin{corollary}\label{cor:shape}
Over a fixed nonzero commutative ring $A$, two diagonal simplicial
tensor modules are isomorphic if and only if one shape is obtained from
the other by permuting its entries and adding the same integer to every entry.
\end{corollary}
\begin{proof}
An isomorphism gives equality of the ranks of the finite free
degree-$p$ modules for every $p\geq0$. Hence their integer
polynomials $\prod_a(x+c_a+1)$ agree at all nonnegative integers,
and therefore agree identically. Their roots, with multiplicities,
show that each offset occurs equally often in both shapes. This is
equivalent to the stated relation between the shapes. Conversely, a permutation
matching the offsets induces a simplicial isomorphism by permuting
coordinates.
\end{proof}

Adding the same integer to every shape entry changes the degree in which
the original tensor occurs, while preserving the simplicial module.

\section{Equivariant contraction}
Let $G=\Stab(\vec s)=\{\sigma\in S_k:n_{\sigma(a)}=n_a\text{ for all }a\}$.
It acts on $X_p$ by permuting tensor axes:
$\sigma E_m=E_{\sigma m}$, where $(\sigma m)_a=m_{\sigma^{-1}(a)}$.
Since the index maps apply the same operation in every coordinate,
this action commutes with all faces and degeneracies. In particular,
it acts on the chain complex $(X_\bullet,\partial)$.
A \emph{$G$-equivariant chain contraction} is a family of $A$-linear maps
$H_p:X_p\to X_{p+1}$ such that
\[
 \partial_{p+1}H_p+H_{p-1}\partial_p=1_{X_p},\qquad
 H_p\sigma=\sigma H_p,
\]
where $X_{-1}=0$ and $H_{-1}=0$.

\begin{proposition}\label{prop:global}
The maps $H_p(E_m)=E_{m+\mathbf1}$ define a $G$-equivariant chain
contraction of $X_\bullet$.
\end{proposition}
\begin{proof}
The shift is well-defined because $m_a\leq p+c_a$ implies
$m_a+1\leq p+1+c_a$. No coordinate of $m+\mathbf1$ is zero,
so deleting the gap at zero gives
$d_0H_p(E_m)=E_m$.

For $p\geq1$ and $1\leq i\leq p+1$, we claim
\begin{equation}\label{eq:shift-faces}
 d_iH_p=H_{p-1}d_{i-1}.
\end{equation}
If $i-1\in\im(m)$, then $i\in\im(m+\mathbf1)$; both sides
vanish on $E_m$ by \eqref{eq:face}. Otherwise, both faces survive,
and the coordinate gap maps satisfy
\[
 \Delta_i^{-1}(m+\mathbf1)
   =\Delta_{i-1}^{-1}(m)+\mathbf1.
\]
Indeed, deleting $i$ subtracts one from a coordinate greater than $i$,
and leaves a coordinate less than $i$ unchanged. Applying this rule
before or after the shift gives the displayed identity, proving
\eqref{eq:shift-faces}. Hence, for $p\geq1$,
\[
 \partial_{p+1}H_p
 =1+\sum_{i=1}^{p+1}(-1)^iH_{p-1}d_{i-1}
 =1-H_{p-1}\partial_p.
\]
In degree zero, some $c_a$ is zero, so every $m\in I_0$ contains
zero. Thus $m+\mathbf1$ contains $1$, giving $d_1H_0=0$ and
$\partial_1H_0=(d_0-d_1)H_0=1$.

Finally, coordinate permutations commute with the shift:
$\sigma(m+\mathbf1)=\sigma m+\mathbf1$. Therefore
$\sigma H_p(E_m)=E_{\sigma m+\mathbf1}=H_p\sigma(E_m)$.
\end{proof}

For constant shapes, $G=S_k$. Define the symmetric submodule by
$\sigma T=T$ for every $\sigma\in S_k$, and the alternating submodule
by $\sigma T=\operatorname{sgn}(\sigma)T$.
The structure maps preserve these submodules because they commute with
$G$. The contraction preserves them as well: if
$\sigma T=\chi(\sigma)T$, where $\chi$ is the trivial or sign character,
then
\[
 \sigma H_p(T)=H_p(\sigma T)=\chi(\sigma)H_p(T).
\]
Restricting the contraction identity therefore contracts both chain
subcomplexes over any $A$; no averaging or division by $k!$ is required.
The normalized contraction $h$ is separately equivariant: its defining
condition $0\in\im(m)$ is invariant under coordinate permutations,
and its nonzero values use the same coordinatewise shift.

Over $\mathbb C$, complex conjugation commutes with the structure maps
and with $H$ and $h$, since their coordinate formulas have real
coefficients. For $k=2$ and a square shape, the axis transposition is
matrix transpose in each degree. Consequently conjugate transpose
$T\mapsto T^*=\overline T^{\mathsf T}$ acts on the underlying real
complexes $X_\bullet$ and $N_\bullet(X)$ and commutes with $H$ and $h$,
respectively. Its $+1$ and $-1$ eigenspaces, the Hermitian and
skew-Hermitian matrices, are real chain subcomplexes on which the
respective contractions restrict.

\begin{remark}
The decreasing filtration
\[
 \mathcal F^tX_p=
 \operatorname{span}_A\{E_m:\min_a m_a\geq p+t\},\qquad t\in\mathbb Z,
\]
is preserved by $\partial$ and $H$. For a surviving face,
$\min_a(\Delta_i^{-1}m)_a\geq\min_a m_a-1$, which gives
$d_i\mathcal F^tX_p\subseteq\mathcal F^tX_{p-1}$.
The shift raises the minimum by exactly one, so
$H_p\mathcal F^tX_p\subseteq\mathcal F^tX_{p+1}$.
The filtration is finite in each degree:
$\mathcal F^{-p}X_p=X_p$ and $\mathcal F^1X_p=0$, since
$\min_a m_a\leq p$. Thus the contraction identity descends to the
associated graded complex. Its homology, the $E_1$ page of the
associated spectral sequence, is zero.
\end{remark}

\section{Generated submodules}
\label{sec:generated}
Throughout this section, $K$ is a field and $X=X(\vec s;K)$.
We describe the submodule $\langle T\rangle$ from Section~2
by its normalized kernel. This gives finite-dimensional data for
computing its homology.

Let $P=K[\Delta^n]$ be the free simplicial $K$-module on the standard
$n$-simplex, so $P_p$ has basis the nondecreasing maps $[p]\to[n]$.
For $T\in X_n$, its representing property gives
\[
 f_T:P_\bullet\longrightarrow X_\bullet,
 \qquad f_T(\alpha)=X(\alpha)(T)
 \quad\bigl(\alpha:[p]\to[n]\bigr).
\]
Its image is the simplicial submodule $\langle T\rangle$ generated by $T$,
and $P/\ker f_T\cong\langle T\rangle$. Put
\[
 C=N(P),\qquad g_T=N(f_T),\qquad L(T)=\ker g_T.
\]
Normalization is exact, so
\begin{equation}\label{eq:generated-normalization}
 N(\langle T\rangle)\cong C/L(T).
\end{equation}
Here $C_p$ has dimension $\binom{n+1}{p+1}$ for $0\leq p\leq n$
and is zero otherwise. Hence the finitely many spaces $L_p(T)$ determine
the entire simplicial kernel by Dold--Kan. They classify the quotient
with its distinguished generating $n$-simplex; they need not determine
its inclusion in $X$, and different kernels can yield isomorphic
simplicial modules after the generator is forgotten.

Homology is computed from normalized complexes. Since $H_0(C)=K$ and
$H_j(C)=0$ for $j>0$, the short exact sequence
$0\to L(T)\to C\to C/L(T)\to0$ gives
\begin{equation}\label{eq:generated-homology}
 \begin{gathered}
 H_j(\langle T\rangle)\cong H_{j-1}(L(T))\qquad(j\geq2),\\
 0\longrightarrow H_1(\langle T\rangle)\longrightarrow H_0(L(T))
 \xrightarrow{\epsilon}K\longrightarrow H_0(\langle T\rangle)
 \longrightarrow0,
 \end{gathered}
\end{equation}
where $\epsilon$ is induced by the augmentation sending every vertex of
$\Delta^n$ to $1$.

For the matrix in Example~\ref{ex:matrix-cycle},
$N(\langle E_{01}\rangle)$ is $K$ in degree $1$ and zero elsewhere:
all its higher simplices are degenerate.

\appendix
\section{Kernel families over an infinite field}
\label{app:kernel-families}
Assume now that $K$ is infinite. Retain the notation $P$, $C$, $g_T$,
and $L(T)$ from Section~\ref{sec:generated}. We parameterize the kernels
on the locus where their dimensions are minimal.

For a finite-dimensional $K$-vector space $V$ and $0\leq r\leq\dim V$,
write $\operatorname{Gr}(r,V)$ for the Grassmannian of $r$-dimensional
linear subspaces of $V$. This is a projective $K$-variety whose points
over a field extension $F/K$ are the $r$-dimensional $F$-linear subspaces
of $V\otimes_K F$. Its tautological subbundle of the trivial bundle
with fiber $V$ has fiber $L$ at the point representing $L\subseteq V$.

View $W=X_n$ as an affine space. For $0\leq p\leq n$, let $r_p$ be
the maximal rank of $g_{T,p}$ and put $\ell_p=\dim C_p-r_p$.
The locus $U\subset W$ on which all these ranks are attained is a
nonempty open subset. Constant-rank kernels define a morphism
\begin{equation}\label{eq:kernel-map}
 \Psi:U\longrightarrow G=\prod_{p=0}^n\operatorname{Gr}(\ell_p,C_p),
 \qquad T\longmapsto(L_p(T))_p.
\end{equation}
Its values satisfy $d_0L_p\subseteq L_{p-1}$, so they are chain
subcomplexes of $C$. By Chevalley's theorem \cite[Tag 054K]{Stacks},
the image $\mathcal M=\Psi(U)$ is constructible. Give its Zariski
closure $Y=\overline{\mathcal M}$ the reduced
structure. When $\mathcal M$ is locally closed, it too is given the
reduced induced structure. Thus $\mathcal M$
records marked quotients of $P$, rather than isomorphism classes of
unmarked simplicial modules. A $K$-variety is \emph{rational over $K$}
if it is birational over $K$ to an affine space $\mathbb A_K^d$.

\begin{proposition}\label{prop:kernel-rationality}
The variety $\overline{\mathcal M}$ is rational over $K$.
\end{proposition}
\begin{proof}
The closure is integral because $U$ is irreducible. Each $g_{T,p}$
depends linearly on $T$. For a tuple $L=(L_p)$, the conditions
\[
 g_{T,p}(L_p)=0\quad(0\leq p\leq n)
\]
therefore define a linear subspace $V_L\subset W$.
Let $\mathcal E\subset W\times Y$ be the incidence scheme cut out by
these conditions on the tautological subbundles $L_p$ pulled back from
the factors of $G$. Over $U\times Y$,
$L_p$ and $\ker g_{T,p}$ are subbundles of equal rank, so inclusion
forces equality. Hence $\mathcal E\cap(U\times Y)$ is the graph of
$\Psi$ as a scheme, and its fiber over $L$ is $U\cap V_L$.

Let $F=K(Y)$ be the function field of $Y$, let $\eta$ be its generic
point, and write $W_F,U_F$ for base change from $K$ to $F$.
Let $V_\eta\subset W_F$ be the linear subspace defined by the incidence
conditions at $\eta$. Its intersection with $U_F$
is nonempty, since $\Psi$ is dominant onto its image closure. Put
$d=\dim_F V_\eta$. Choose $d$ coordinate functions on $W$ whose
restriction gives an isomorphism $V_\eta\to F^d$, and denote their
joint map by $\pi:W\to K^d$. Solving the defining linear equations
expresses the unique element of $V_\eta$ with $\pi(T)=c$ rationally
in the kernel coordinates and linearly in $c$. The condition that this
element belongs to $U_F$ defines a nonempty open subset of $F^d$.
Since $K$ is infinite, $K^d$ is Zariski dense in $F^d$; choose
$c\in K^d$ in that open subset.

After restricting to a nonempty open subset of
$\overline{\mathcal M}$, the subspaces $V_L$ have constant dimension,
$\pi|_{V_L}$ is invertible, and the solutions with $\pi(T)=c$
belong to $U$. Every fiber there consequently meets the affine slice
$S=\pi^{-1}(c)$ in exactly one point, varying rationally with $L$.
These solutions and $\Psi|_S$ are inverse on nonempty open subsets.
As $S\cong\mathbb A_K^{\dim W-d}$, the closure is rational.
\end{proof}

For $V=K^3$, a plane is the kernel of a nonzero linear form $(a,b,c)$,
unique up to a nonzero scalar. Thus
\[
 \operatorname{Gr}(2,K^3)\xrightarrow{\ \sim\ }\mathbb P_K^2,
 \qquad \ker(a,b,c)\longmapsto[a:b:c],
\]
where $[a:b:c]$ denotes the corresponding homogeneous coordinates.

\begin{proposition}\label{prop:three-by-three}
For shape $(3,3)$, the generic normalized kernel dimensions are
$(\ell_0,\ell_1,\ell_2)=(2,0,0)$, and the kernel image is
\[
 \mathcal M\cong\mathbb P_K^2\setminus\{[1:1:1]\}.
\]
For every $T\in U$,
\[
 H_1(\langle T\rangle)\cong K,
 \qquad H_j(\langle T\rangle)=0\quad(j\ne1).
\]
\end{proposition}
\begin{proof}
Write $T=(t_{ij})$ and $(a,b,c)=(t_{00},t_{11},t_{22})$.
In degree zero, $f_T$ is the row $(a,b,c):K^3\to K$.
In degree one, the columns corresponding to the three constant maps
$[1]\to[2]$ are $aJ,bJ,cJ$, where $J$ is the all-ones $2\times2$
matrix. When $(a,b,c)\ne0$, they span $KJ$. Modulo this line, the
other three columns are represented by
\[
 B_T=\begin{pmatrix}
 t_{01}-a&t_{02}-a&t_{12}-b\\
 t_{10}-a&t_{20}-a&t_{21}-b\\
 b-a&c-a&c-b
 \end{pmatrix}.
\]
Its first two rows are arbitrary, while its last row is nonzero
exactly when $a,b,c$ are not all equal. Consequently $f_{T,1}$ can
have rank $4$ for a specified diagonal precisely when that diagonal
is nonconstant. Such a choice gives $\rank g_{T,0}=1$ and
$\rank g_{T,1}=3$, using the natural decomposition
$P_1=C_1\oplus s_0C_0$ and its counterpart for $X$.

The normalized simplex complex $C=N(K[\Delta^2])$ has dimensions
$(3,3,1)$ and injective differential $C_2\to C_1$. Since $g_{T,1}$
is injective, the chain-map identity makes $g_{T,2}$ injective as
well. These are maximal ranks in all degrees, and
$U=\{T:\det B_T\ne0\}$. Thus on $U$ the kernel
complex is $\ker(a,b,c)$ in degree zero and zero elsewhere.
Conversely, every nonconstant projective diagonal is realized by
choosing the first two rows of the displayed matrix to complete its
last row to a basis. This construction works over every residue field.
Under the displayed identification of $\operatorname{Gr}(2,K^3)$ with
$\mathbb P_K^2$, the kernel map therefore has exactly the stated image.

The restriction of the augmentation $(1,1,1)$ to $\ker(a,b,c)$ is
nonzero, because $(a,b,c)$ is not proportional to $(1,1,1)$. It is
therefore surjective with one-dimensional kernel. Formula
\eqref{eq:generated-homology} proves the homology assertions.
\end{proof}

For comparison, the all-ones $3\times3$ tensor generates the constant
simplicial module $K$: every simplicial operator sends it to an
all-ones tensor. Its homology is $K$ in degree zero and zero in
positive degrees. Thus generated submodules distinguish tensors by
homology even when their common ambient module is contractible.


\bibliographystyle{plain}
\bibliography{dstm_short}
\end{document}
